\documentclass[11pt]{article}
\usepackage[margin=1in]{geometry}
\usepackage{amsmath,amssymb,amsthm}
\usepackage{booktabs}
\usepackage{graphicx}
\usepackage{hyperref}
\usepackage{xcolor}
\usepackage{microtype}
\usepackage{natbib}

\newtheorem{proposition}{Proposition}
\newtheorem{remark}{Remark}

\title{Fisher--Rao Loss Does Not Tolerate Symmetric Label Noise:\\
A Comparative Study of Distribution Divergences for Noisy-Label Learning}

\author{Maruan Ottoni}

\date{\today}

\begin{document}
\maketitle

\begin{abstract}
Learning under noisy labels is a central challenge in modern machine learning.
While noise-tolerant losses such as Mean Absolute Error (MAE) and Generalized
Cross Entropy (GCE) are well-understood theoretically, the behavior of
information-geometric divergences such as the Fisher--Rao (FR) distance has
not been systematically evaluated in this setting.
We conduct a controlled benchmark comparing six objectives on the UCI Digits
dataset across five symmetric noise rates (0--80\%) and one asymmetric noise
rate (40\%).
Our main finding is negative but informative: the FR loss and the closely
related Hellinger distance both \emph{under-perform} standard cross-entropy
(KL) at moderate symmetric noise rates (sym\,40: $\Delta\mathrm{acc} = -3.2\%$,
$p=0.002$, 0/10 paired wins; sym\,60: $-2.0\%$, $p=0.004$, 0/10 wins).
By contrast, MAE achieves $+19.6\%$ at sym\,40 ($p=0.002$, 10/10 wins) and
GCE/SCE achieve $+6.1$--$6.3\%$.
We trace this gap to the Ghosh et al.\ noise-tolerance condition: a loss is
Bayes-optimal under symmetric noise if and only if its per-class values sum to
a constant, a property that MAE satisfies exactly but FR does not.
These results sharpen the understanding of where FR-based training is and is
not beneficial.
\end{abstract}

% ------------------------------------------------------------------
\section{Introduction}
% ------------------------------------------------------------------

Label noise is pervasive in large-scale datasets assembled from web crawls,
crowd-sourcing, and automated pipelines.
Even modest noise rates can substantially degrade test accuracy, and the
classical cross-entropy (KL) loss is known to be brittle in this setting.

Two families of remedies have been proposed.
\emph{Noise-tolerant losses} replace cross-entropy with objectives whose
gradients are inherently bounded or whose risk minimizers are provably
Bayes-optimal under symmetric noise~\citep{ghosh2017robust,
zhang2018generalized, wang2019symmetric}.
\emph{Robust training procedures} apply data selection, teacher--student
distillation, or sample weighting to reduce the influence of noisy
examples~\citep{han2018coteaching, li2020dividemix}.

A complementary line of work draws on \emph{information geometry}: replacing
KL divergence with the Fisher--Rao geodesic distance or the Hellinger distance
in the training objective.
These divergences are bounded, symmetric, and properly metric on the simplex,
properties that are absent from KL and that have been argued to confer
robustness~\citep{bottoni2026fisherrao}.
However, no systematic empirical study has compared FR/Hellinger against
established noise-tolerant baselines across a range of noise conditions.

We fill this gap with a controlled experiment whose design is simple enough to
be fully reproducible on a single laptop.
Our contributions are:
\begin{enumerate}
  \item A six-way comparison of objectives (KL, Fisher--Rao, Hellinger, GCE,
    MAE, SCE) across six label-noise regimes on UCI Digits.
  \item A theoretical explanation via the Ghosh et al.\ condition showing why
    FR/Hellinger cannot be noise-tolerant.
  \item Practical guidance: use MAE or GCE for symmetric noise; reserve
    Fisher--Rao for settings with structured overconfident soft-target noise.
\end{enumerate}

% ------------------------------------------------------------------
\section{Background}
% ------------------------------------------------------------------

\subsection{Label noise types}

Let $(x, y) \sim \mathcal{D}$ be a clean example and $\tilde{y}$ the observed
noisy label.

\paragraph{Symmetric noise at rate $\varepsilon$.}
Each label is replaced uniformly at random by another class with probability
$\varepsilon$:
\[
  \tilde{y} \;=\;
  \begin{cases}
    y & \text{with probability } 1-\varepsilon \\
    \mathrm{Uniform}(\mathcal{Y} \setminus \{y\}) & \text{with probability } \varepsilon.
  \end{cases}
\]

\paragraph{Asymmetric (class-conditional) noise at rate $\varepsilon$.}
Each class $c$ is mapped to class $(c+1)\!\mod C$ with probability $\varepsilon$
and remains $c$ otherwise.

\subsection{Objectives}

All objectives receive a one-hot target $\mathbf{y}$ (possibly corrupted) and
a predicted probability vector $\mathbf{p} = \mathrm{softmax}(\mathbf{z})$.

\paragraph{Cross-entropy (KL).}
$\mathcal{L}_\mathrm{KL} = -\sum_k y_k \log p_k$.

\paragraph{Fisher--Rao (FR).}
The geodesic distance on the categorical simplex, squared:
$\mathcal{L}_\mathrm{FR} = 4\arccos^2\!\bigl(\sum_k \sqrt{y_k p_k}\bigr)$.

\paragraph{Hellinger.}
$\mathcal{L}_\mathrm{H} = \tfrac{1}{2}\sum_k (\sqrt{y_k} - \sqrt{p_k})^2$.

\paragraph{Generalized Cross Entropy (GCE)~\citep{zhang2018generalized}.}
$\mathcal{L}_\mathrm{GCE} = (1 - p_{\tilde{y}}^q)/q$ with $q=0.7$.

\paragraph{Mean Absolute Error (MAE)~\citep{ghosh2017robust}.}
$\mathcal{L}_\mathrm{MAE} = \sum_k |y_k - p_k|$.

\paragraph{Symmetric Cross Entropy (SCE)~\citep{wang2019symmetric}.}
$\mathcal{L}_\mathrm{SCE} = 0.1\cdot\mathcal{L}_\mathrm{KL} +
1.0\cdot\mathcal{L}_\mathrm{RCE}$, where $\mathcal{L}_\mathrm{RCE} = -\sum_k
p_k \log y_k$ is the reverse KL.

\subsection{Noise-tolerance condition}

\begin{proposition}[Ghosh et al.\ 2017, Theorem~1]
  \label{prop:ghosh}
  A loss $\ell(\mathbf{f}(x), \mathbf{e}_k)$ is Bayes-optimal under symmetric
  label noise with rate $\varepsilon < (C-1)/C$ if and only if
  $\sum_{k=1}^C \ell(\mathbf{f}(x), \mathbf{e}_k) = \mathrm{const}$
  for all $\mathbf{f}(x) \in \Delta^{C-1}$.
\end{proposition}

MAE satisfies this condition: $\sum_k \|\mathbf{p} - \mathbf{e}_k\|_1 = 2(C-1)$
for any $\mathbf{p}$~\citep{ghosh2017robust}.
Cross-entropy does not ($\sum_k -\log p_k$ is unbounded).

\begin{remark}
  The Fisher--Rao loss $4\arccos^2(\sqrt{p_k})$ sums to
  $4\sum_k \arccos^2(\sqrt{p_k})$, which depends on $\mathbf{p}$ and is
  therefore \emph{not} constant.
  Similarly, $\sum_k \tfrac{1}{2}(\sqrt{p_k} - 1)^2 = \tfrac{1}{2}(1 -
  2\sum_k \sqrt{p_k} + C)$, which depends on $\sum_k\sqrt{p_k}$ and is
  not constant.
  Neither FR nor Hellinger satisfies the noise-tolerance condition.
\end{remark}

% ------------------------------------------------------------------
\section{Experimental Setup}
% ------------------------------------------------------------------

\paragraph{Dataset.}
UCI Handwritten Digits (1,797 samples, 64 features, 10 classes).
80/20 train/test split stratified by class; features standardized.

\paragraph{Model.}
Two-hidden-layer MLP with 256 units, batch normalisation, and ReLU activations.

\paragraph{Training.}
Adam optimiser, learning rate $10^{-3}$, weight decay $10^{-4}$, cosine
annealing over 100 epochs, batch size 128, gradient clipping at 5.0.

\paragraph{Noise injection.}
Noise is applied to training labels only; the test set is always clean.
For symmetric noise at rate $\varepsilon$, we sample $\lfloor \varepsilon N
\rfloor$ training indices uniformly at random and replace their labels with a
uniformly random alternative class.
For asymmetric noise at rate $\varepsilon = 0.4$, each label $c$ is set to
$(c+1)\!\mod 10$ independently with probability $0.4$.

\paragraph{Evaluation.}
Ten independent seeds per (noise regime, objective) cell.
Primary metric: test accuracy.
Secondary metrics: ECE (expected calibration error), Brier score, NLL.
Statistical comparison: two-sided Wilcoxon signed-rank test against KL with
paired seeds; oriented gain $= \mathrm{mean}(\Delta\mathrm{acc})$ with positive
meaning improvement over KL.

% ------------------------------------------------------------------
\section{Results}
% ------------------------------------------------------------------

\subsection{Accuracy under symmetric noise}

Table~\ref{tab:accuracy} reports mean test accuracy ($\pm$ std across 10 seeds)
for each (noise regime, objective) cell.

\begin{table}[t]
  \centering
  \caption{Mean test accuracy (\%) on UCI Digits.
    Bold: best non-KL objective per noise regime.
    Underline: statistically worse than KL ($p<0.01$, Wilcoxon).}
  \label{tab:accuracy}
  \small
  \begin{tabular}{lcccccc}
    \toprule
    Noise & KL (CE) & FR & Hellinger & GCE & MAE & SCE \\
    \midrule
    Clean    & 97.9 & 97.9 & 97.9 & 98.0 & 98.0 & 97.9 \\
    Sym 20\% & 86.2 & 86.4 & 86.5 & \textbf{92.4} & \textbf{96.6} & 91.9 \\
    Sym 40\% & 70.0 & \underline{66.8} & \underline{67.2} & \textbf{76.3} & \textbf{89.6} & \textbf{76.1} \\
    Sym 60\% & 47.1 & \underline{45.1} & \underline{45.4} & \textbf{50.7} & \textbf{68.2} & \textbf{50.6} \\
    Sym 80\% & 21.7 & 22.0 & 21.9 & 23.1 & \textbf{29.5} & 23.6 \\
    Asym 40\% & 61.1 & 60.3 & 60.7 & 60.9 & \textbf{67.3} & 61.2 \\
    \bottomrule
  \end{tabular}
\end{table}

Key observations:

\begin{enumerate}
  \item \textbf{MAE dominates at all symmetric noise levels.}
    At sym\,40, MAE achieves $89.6\%$ vs KL's $70.0\%$ ($+19.6\%$,
    $p=0.002$, 10/10 paired wins).
    At sym\,60, the gap widens to $+21.1\%$ ($p=0.002$, 10/10 wins).
  \item \textbf{GCE and SCE are competitive.}
    Both achieve $+6.1$--$6.3\%$ over KL at sym\,40 ($p=0.002$, 10/10 wins).
  \item \textbf{FR and Hellinger are significantly worse than KL at sym\,40
    and sym\,60.}
    FR: $-3.2\%$ at sym\,40 ($p=0.002$, 0/10 wins);
    $-2.0\%$ at sym\,60 ($p=0.004$, 0/10 wins).
    Hellinger shows the same pattern ($-2.9\%$ at sym\,40, $p=0.004$, 1/10
    wins).
  \item \textbf{At sym\,80, differences vanish.}
    At very high noise the signal-to-noise ratio collapses; all objectives
    converge near random ($\approx 22$--$29\%$).
  \item \textbf{Asymmetric noise: only MAE wins significantly.}
    $+6.2\%$ ($p=0.002$, 10/10 wins).
    FR is not significantly different from KL ($p=0.057$, 2/10 wins).
\end{enumerate}

\subsection{Calibration and Brier score}

Across all symmetric noise regimes, MAE consistently achieves the lowest Brier
score and NLL, consistent with its theoretical guarantees.
GCE and SCE track closely.
FR and Hellinger show elevated Brier scores and NLL at sym\,40--60, confirming
that the accuracy drop reflects a systematic failure rather than a
variance artefact.

% ------------------------------------------------------------------
\section{Discussion}
% ------------------------------------------------------------------

\paragraph{Why does FR fail under symmetric noise?}
The failure is explained by Proposition~\ref{prop:ghosh}: FR's per-class loss
values sum to a quantity that depends on the predicted distribution, so the
risk minimizer under symmetric noise is \emph{not} the Bayes classifier.
In contrast, MAE's constant-sum property makes it immune to symmetric noise by
construction.

The situation is conceptually similar to the KL vs.\ MAE comparison: KL also
lacks the constant-sum property and is therefore brittle to symmetric noise.
But KL at least benefits from numerical stability and well-understood
optimisation dynamics.
FR combines the worst of both worlds: it lacks the constant-sum property
\emph{and} its gradient vanishes near one-hot distributions (where
$\arccos(\sqrt{p_k}) \approx 0$), leaving the optimiser with weak signal
precisely when training examples are corrupted.

\paragraph{When does FR help?}
Previous work~\citep{bottoni2026fisherrao} shows FR outperforming KL in the
\emph{soft-label} setting, where targets are themselves probability distributions
rather than one-hot vectors.
When the target distribution is genuinely diffuse (e.g.\ knowledge distillation
or soft crowdsourced annotations), FR's boundedness prevents overconfident
responses to mislabelled soft targets.
The present experiment uses one-hot targets corrupted by a random noise process,
which is a fundamentally different regime.

\paragraph{Practical recommendation.}
For symmetric label noise, use MAE (strongest guarantee, best empirical
performance) or GCE/SCE (good performance with KL-like optimisation dynamics).
For structured noise in probability targets (overconfident wrong distributions),
the Fisher--Rao loss may be preferable.

% ------------------------------------------------------------------
\section{Related Work}
% ------------------------------------------------------------------

\citet{ghosh2017robust} established the noise-tolerance condition and showed MAE
satisfies it.
\citet{zhang2018generalized} proposed GCE as a differentiable interpolation
between CE and MAE.
\citet{wang2019symmetric} introduced SCE to combine forward and reverse KL for
label noise robustness.
\citet{ma2020normalized} showed that normalised versions of many losses
(including MAE) satisfy the noise-tolerance condition.

On the information-geometry side, \citet{bottoni2026fisherrao} introduced the
Fisher--Rao loss for distribution comparison in the soft-label and
dimensionality-reduction settings, establishing theoretical advantages over KL
in that context.
The present work is the first to benchmark FR specifically under the
symmetric/asymmetric label noise model of \citet{ghosh2017robust}.

% ------------------------------------------------------------------
\section{Limitations}
% ------------------------------------------------------------------

\begin{itemize}
  \item Experiments use UCI Digits (1,797 samples, 64 features).
    Results on larger datasets with convolutional architectures (e.g.\ CIFAR-10,
    Clothing-1M) may differ.
  \item We evaluate only synthetic noise.
    Real-world noise has complex class-dependent structure not captured by
    the symmetric or single-hop asymmetric models.
  \item We do not evaluate co-learning, teacher--student, or semi-supervised
    approaches; the comparison is restricted to the choice of loss function alone.
  \item Ten seeds gives a minimum achievable $p$-value of $\approx 0.002$
    (Wilcoxon on $n=10$ pairs); very small effects cannot be detected reliably.
\end{itemize}

% ------------------------------------------------------------------
\section{Conclusion}
% ------------------------------------------------------------------

We conducted a systematic evaluation of the Fisher--Rao geodesic distance as a
training objective under label noise.
Contrary to the intuition that a bounded, symmetric divergence should be more
robust, FR under-performs standard cross-entropy at symmetric noise rates of
40--60\% and shows no advantage at other tested rates.
The theoretical explanation is precise: FR does not satisfy the Ghosh et al.\
noise-tolerance condition, so its risk minimizer under symmetric noise is not
the Bayes classifier.

MAE remains the theoretically principled and empirically strongest choice for
symmetric label noise, achieving up to $+21\%$ accuracy over KL at 60\% noise.
GCE and SCE offer a practical compromise between noise robustness and
optimisation stability.

The Fisher--Rao loss retains its advantages in settings involving soft
probability targets, and future work should explore hybrid or adaptive
objectives that switch between FR and MAE depending on the estimated noise
structure.

\bibliographystyle{plainnat}
\bibliography{references}

\end{document}
